\documentclass[UTF8,aspectratio=169,10pt]{ctexbeamer}

\usetheme{Madrid}
\usecolortheme{default}
\usepackage{fontspec}
\usepackage{booktabs}
\usepackage{array}
\usepackage{multirow}
\usepackage{tikz}
\usepackage{xcolor}
\usepackage{listings}
\usetikzlibrary{arrows.meta,positioning,fit,calc,shapes.geometric}

\setCJKmainfont{Microsoft YaHei}
\setmainfont{Arial}
\hypersetup{colorlinks=true,urlcolor=blue}

\definecolor{memoBlue}{RGB}{37,99,235}
\definecolor{memoGreen}{RGB}{22,101,52}
\definecolor{memoOrange}{RGB}{180,83,9}
\definecolor{memoPurple}{RGB}{109,40,217}
\definecolor{memoGray}{RGB}{71,85,105}
\definecolor{memoLight}{RGB}{248,250,252}

\setbeamercolor{title}{fg=white}
\setbeamercolor{subtitle}{fg=white}
\setbeamercolor{author}{fg=white}
\setbeamercolor{date}{fg=white}
\setbeamercolor{frametitle}{fg=white,bg=memoBlue}
\setbeamercolor{block title}{fg=white,bg=memoBlue}
\setbeamercolor{block body}{fg=black,bg=memoLight}
\setbeamercolor{alerted text}{fg=memoOrange}

\setbeamertemplate{navigation symbols}{}
\setbeamertemplate{footline}[frame number]

\lstdefinestyle{memo}{
  basicstyle=\ttfamily\scriptsize,
  breaklines=true,
  columns=fullflexible,
  keepspaces=true,
  showstringspaces=false,
  frame=single,
  rulecolor=\color{memoGray!40},
  backgroundcolor=\color{memoLight}
}
\lstset{style=memo}

\newcommand{\code}[1]{\texttt{\small #1}}
\newcommand{\smallcode}[1]{\texttt{\scriptsize #1}}
\newcommand{\metric}[2]{\textbf{\textcolor{memoBlue}{#1}}{\scriptsize\;#2}}
\newcommand{\sourcefoot}[1]{%
  \vfill
  {\scriptsize\textcolor{memoGray}{Source: }%
  \parbox[t]{0.84\textwidth}{\raggedright\ttfamily\scriptsize #1}}%
}
\newcommand{\explainfoot}[1]{%
  \vspace{0.35em}
  {\scriptsize\color{memoGray}#1}
}

\title[MEMO-Appflow]{MEMO-Appflow}
\subtitle{从 Android OS 证据到端侧 MAPLE 推理，再回到应用推荐与系统调度}
\author{讲解人：Jingyi Guo}
\date{2026-07-02}

\begin{document}

{
\setbeamercolor{background canvas}{bg=memoBlue}
\begin{frame}[plain]
  \vspace{2.0em}
  \begin{center}
    {\Huge\bfseries\textcolor{white}{MEMO-Appflow}}\\[1.2em]
    {\Large\textcolor{white}{从 Android OS 证据到端侧 MAPLE 推理}}\\[0.4em]
    {\Large\textcolor{white}{再回到应用推荐与系统调度}}\\[2.0em]
    {\normalsize\textcolor{white}{讲解人：Jingyi Guo}}\\[0.5em]
    {\small\textcolor{white}{组员：Mosha Huang \quad Kewei Chen \quad Chengyu Fan}}\\[0.8em]
    {\small\textcolor{white}{rooted Pixel 5 \quad | \quad raw eBPF/libbpf \quad | \quad local MAPLE}}\\[0.3em]
    {\small\textcolor{white}{Top-3 真实应用推荐 + scheduler bits 系统调度}}
  \end{center}
\end{frame}
}

\begin{frame}{汇报路线}
  \tableofcontents
\end{frame}

\section{从 app + OS 到 app + OS}

\begin{frame}{项目不是单纯的 app 推荐}
  \begin{columns}[T,totalwidth=\textwidth]
    \begin{column}{0.48\textwidth}
      \begin{block}{用户看到的现象}
        \begin{itemize}
          \item 刚打开聊天软件，马上想拍照或分享。
          \item 正在浏览网页，网络和渲染压力升高。
          \item 手机缓存很多 app 后，启动等待变长。
        \end{itemize}
      \end{block}
    \end{column}
    \begin{column}{0.48\textwidth}
      \begin{block}{OS 真实承受的压力}
        \begin{itemize}
          \item Binder 服务调用、前台切换、进程 exec。
          \item UDP send/recv、文件 open、runtime loading。
          \item vmscan reclaim、SurfaceFlinger、调度事件。
        \end{itemize}
      \end{block}
    \end{column}
  \end{columns}
  \vspace{0.8em}
  \begin{alertblock}{MEMO-Appflow 的核心判断}
    只看 app 序列不够，只看内核事件也不够。我们把 app 行为和 OS 证据合并成端侧模型上下文，再让模型同时输出用户接下来可能用的 app 和系统应该执行的调度动作。
  \end{alertblock}
\end{frame}

\begin{frame}{宏观闭环：系统观察、模型决策、系统执行}
  \centering
  \resizebox{0.95\textwidth}{!}{
  \begin{tikzpicture}[
    font=\small,
    box/.style={draw, rounded corners, align=center, minimum width=3.2cm, minimum height=0.9cm},
    wide/.style={draw, rounded corners, align=center, minimum width=4.2cm, minimum height=0.9cm},
    arrow/.style={-{Latex[length=2.2mm]}, thick}
  ]
    \node[wide, fill=blue!9] (os) {Android OS\\tracepoints 内核事件 + system state 系统状态};
    \node[wide, fill=blue!9, right=1.1cm of os] (app) {User/App Layer\\foreground app timeline};
    \node[wide, fill=green!10, below=1.0cm of $(os)!0.5!(app)$] (struct) {Device-side structuring 端侧结构化\\timeline windows 时间窗 + compressed eBPF 压缩信号};
    \node[wide, fill=orange!14, below=1.0cm of struct] (maple) {Local MAPLE reasoning\\Stage 1 category + Stage 2 app id + Stage 3 scheduler bits};
    \node[box, fill=purple!10, below left=1.0cm and -0.2cm of maple] (top3) {Top-3 real apps\\PackageManager 真实应用映射};
    \node[box, fill=purple!10, below right=1.0cm and -0.2cm of maple] (act) {OS actions\\warm launch / trim / refresh};
    \node[wide, fill=gray!12, below=1.0cm of $(top3)!0.5!(act)$] (ui) {Widget + App UI\\可见推荐 + 动作记录};

    \draw[arrow] (os) -- (struct);
    \draw[arrow] (app) -- (struct);
    \draw[arrow] (struct) -- (maple);
    \draw[arrow] (maple) -- (top3);
    \draw[arrow] (maple) -- (act);
    \draw[arrow] (top3) -- (ui);
    \draw[arrow] (act) -- (ui);
  \end{tikzpicture}}
  \vspace{0.25em}
  \scriptsize
  MAPLE 是本项目本地大模型推理模块；scheduler bits 是模型输出的一串 0/1 系统动作开关；
  Widget 是桌面小组件，用来展示可点击的 Top-3 app。这条链路在手机内部运行，host 只负责安装、推模型、拉报告。
\end{frame}

\begin{frame}{手机内运行边界}
  \begin{columns}[T,totalwidth=\textwidth]
    \begin{column}{0.48\textwidth}
      \begin{block}{Phone-local（手机本地执行）}
        \begin{itemize}
          \item raw eBPF/libbpf collector：原生内核事件采集器
          \item foreground app sampler
          \item system state collector：内存、网络、电量等状态采集器
          \item Kotlin parser and scenario builder：Kotlin 解析与场景构造
          \item local MAPLE native inference：本地 MAPLE 推理
          \item Top-3 mapping and ActionExecutor：Top-3 映射与动作执行
          \item desktop Widget update：桌面小组件刷新
        \end{itemize}
      \end{block}
    \end{column}
    \begin{column}{0.48\textwidth}
      \begin{block}{Host only}
        \begin{itemize}
          \item install APK
          \item push MAPLE binary and GGUF model
          \item trigger experiment during debugging
          \item pull reports for inspection
          \item compile LaTeX / write paper
        \end{itemize}
      \end{block}
      \begin{alertblock}{设计原则}
        产品逻辑不能依赖电脑脚本；拔线后手机仍能通过 App 和 Widget 运行核心闭环。
      \end{alertblock}
    \end{column}
  \end{columns}
\end{frame}

\section{eBPF 与 Kernel 采集层}

\begin{frame}{为什么要引入 eBPF}
  {\scriptsize eBPF 是 Linux/Android 内核里的安全观测机制。我们用它在内核事件发生时记录证据，而不是只看 app 层日志。}
  \vspace{0.35em}
  \begin{columns}[T,totalwidth=\textwidth]
    \begin{column}{0.46\textwidth}
      \begin{block}{App 序列能告诉我们}
        \begin{itemize}
          \item 用户刚打开了哪个应用。
          \item 应用切换的大致顺序。
          \item 最近是否处于社交、浏览、媒体等工作流。
        \end{itemize}
      \end{block}
    \end{column}
    \begin{column}{0.50\textwidth}
      \begin{block}{eBPF 能补上}
        \begin{itemize}
          \item 网络收发是否活跃：\code{sendto}/\code{recvfrom}。
          \item 服务调用是否密集：Binder transaction。
          \item 内存是否紧张：direct reclaim / kswapd。
          \item UI 是否繁忙：sched、input、SurfaceFlinger hints。
        \end{itemize}
      \end{block}
    \end{column}
  \end{columns}
  \vspace{0.6em}
  \centering
  \large App sequence describes intent; eBPF describes the hidden OS work behind that intent.
\end{frame}

\begin{frame}{采集的 14 个 tracepoints}
  \scriptsize
  Tracepoint 是内核预留的稳定事件点，例如系统调用进入、Binder transaction、进程 exec。collector attach 到这些点上收真实 OS 证据。
  \vspace{0.35em}
  \begin{tabular}{lll}
    \toprule
    Group & Tracepoint & 用途 \\
    \midrule
    IPC & \code{tracepoint/binder/binder\_transaction} & Android service / Binder pressure \\
    File & \code{syscalls/sys\_enter\_openat} & file I/O, runtime loading \\
    File & \code{syscalls/sys\_enter\_openat2} & file I/O, camera/media file hints \\
    Network & \code{syscalls/sys\_enter\_sendto} & network send, especially UDP \\
    Network & \code{syscalls/sys\_enter\_recvfrom} & network receive, especially UDP \\
    Memory & \code{vmscan/mm\_vmscan\_direct\_reclaim\_begin} & direct reclaim starts \\
    Memory & \code{vmscan/mm\_vmscan\_direct\_reclaim\_end} & direct reclaim ends \\
    Memory & \code{vmscan/mm\_vmscan\_kswapd\_wake} & kswapd wakeup \\
    Scheduling & \code{sched/sched\_switch} & CPU scheduling activity \\
    Scheduling & \code{sched/sched\_wakeup} & wakeup pressure \\
    Process & \code{sched/sched\_process\_fork} & process lifecycle \\
    Process & \code{sched/sched\_process\_exit} & process lifecycle \\
    Process & \code{sched/sched\_process\_exec} & process launch \\
    Input & \code{input/input\_event} & touch/key event \\
    \bottomrule
  \end{tabular}
  \vspace{0.5em}
  \begin{block}{实现位置}
    raw BPF object（BPF 字节码）和 collector（采集器）部署到 \code{/data/local/tmp/memo}；Android 端通过 root shell 启动。
  \end{block}
\end{frame}

\begin{frame}{Pixel 5 上的 kernel/runtime 条件}
  \begin{columns}[T,totalwidth=\textwidth]
    \begin{column}{0.49\textwidth}
      \begin{block}{Device}
        \begin{itemize}
          \item Google Pixel 5
          \item arm64-v8a
          \item Magisk root
          \item custom Android kernel path
          \item SELinux relaxed for tracing/debug deployment
        \end{itemize}
      \end{block}
    \end{column}
    \begin{column}{0.49\textwidth}
      \begin{block}{Kernel features}
        \begin{itemize}
          \item \code{CONFIG\_BPF}
          \item \code{CONFIG\_BPF\_SYSCALL}
          \item \code{CONFIG\_BPF\_JIT}
          \item \code{CONFIG\_FTRACE}
          \item \code{CONFIG\_FTRACE\_SYSCALLS}
          \item \code{CONFIG\_KPROBES}, \code{CONFIG\_UPROBES}
        \end{itemize}
      \end{block}
    \end{column}
  \end{columns}
  \vspace{0.4em}
  \begin{alertblock}{工程选择}
    bpftrace 曾用于探索，但最终产品路径以 raw eBPF/libbpf collector 为主，减少运行时依赖，也更适合 Android 端部署。
  \end{alertblock}
\end{frame}

\begin{frame}{tracepoint 到 counter key}
  \small
  \begin{columns}[T,totalwidth=\textwidth]
    \begin{column}{0.54\textwidth}
      \begin{tabular}{ll}
        \toprule
        tracepoint group & counter key / type \\
        \midrule
        binder transaction & 1 / binder \\
        openat, openat2 & 2 / file \\
        vmscan begin/end/kswapd & 3 / memory \\
        sendto & 4 / network-send \\
        recvfrom & 5 / network-recv \\
        sched switch/wakeup & 6 / sched \\
        process fork & 7 / process-fork \\
        process exit & 8 / process-exit \\
        process exec & 9 / process-exec \\
        input event & 10 / input \\
        \bottomrule
      \end{tabular}
    \end{column}
    \begin{column}{0.42\textwidth}
      \begin{block}{为什么先计数}
        \begin{itemize}
          \item eBPF 原始行可能上万条。
          \item 大量重复事件直接进 prompt 会浪费上下文长度。
          \item counter 保存 count/rate，保留“强度”和“密度”。
          \item raw trace 仍保留在报告目录用于复查。
        \end{itemize}
      \end{block}
    \end{column}
  \end{columns}
  \explainfoot{左边每一行表示一个内核 tracepoint 族会被归成哪一种计数键。比如 \code{sendto/recvfrom} 都是网络收发证据，计数后变成 \code{network-send/network-recv}。这样做不是丢掉 eBPF，而是把重复低层事件压成模型能读懂的“某类系统活动发生了多少次、每秒多密集”。}
\end{frame}

\begin{frame}[fragile]{真实采集窗口长什么样}
  \begin{columns}[T,totalwidth=\textwidth]
    \begin{column}{0.50\textwidth}
      \begin{block}{来自 30-record 实验的窗口摘要}
\begin{lstlisting}
raw_trace_path:
/sdcard/MEMO/logs/realtime_1782914935996.trace

app_samples_path:
/sdcard/MEMO/logs/realtime_1782914935996_apps.tsv

raw_ebpf_events: 13733
app_segments: 6
timeline_windows: 2
\end{lstlisting}
      \end{block}
    \end{column}
    \begin{column}{0.46\textwidth}
      \begin{block}{进入 MAPLE 前会被压缩}
\begin{lstlisting}
{
  "event_type": "MEMO_RECVFROM",
  "detail": "recvfrom",
  "count": 258462,
  "rate_per_sec": 1367.2
}
\end{lstlisting}
      \end{block}
    \end{column}
  \end{columns}
  \explainfoot{这一页展示的是一次手机端真实采集窗口：raw trace 是原始 eBPF 文件，app samples 是同一时间内前台 app 记录。右边的 \code{count} 和 \code{rate\_per\_sec} 表示同类事件在窗口里重复出现的次数和频率，避免把几十万条相似网络收包逐行塞进 prompt。}
  \sourcefoot{docs/real\_device\_experiments/synthetic\_user\_30\_2026\_07\_01}
\end{frame}

\begin{frame}{Android 端采集运行时：代码结构概览}
  \centering
  \resizebox{0.92\textwidth}{!}{
  \begin{tikzpicture}[
    font=\small,
    flowbox/.style={draw, rounded corners, align=center, minimum width=3.1cm, minimum height=0.8cm},
    arrow/.style={-{Latex[length=2mm]}, thick}
  ]
    \node[flowbox, fill=blue!8] (svc) {\code{EBPFCollectorService}\\前台服务入口};
    \node[flowbox, fill=blue!8, right=0.9cm of svc] (deploy) {\code{DeviceCollectorDeployer}\\部署可执行文件};
    \node[flowbox, fill=green!10, right=0.9cm of deploy] (collector) {启动\\\code{memo\_libbpf\_collector}};
    \node[flowbox, fill=green!10, below=0.8cm of collector] (sampler) {启动前台\\app 采样器};
    \node[flowbox, fill=orange!12, left=0.9cm of sampler] (stop) {停止采集器\\读取 MEMO 行};
    \node[flowbox, fill=purple!10, left=0.9cm of stop] (parse) {\code{EBPFTraceParser}\\Kotlin 事件对象};
    \draw[arrow] (svc) -- (deploy);
    \draw[arrow] (deploy) -- (collector);
    \draw[arrow] (collector) -- (sampler);
    \draw[arrow] (sampler) -- (stop);
    \draw[arrow] (stop) -- (parse);
  \end{tikzpicture}}
  \vspace{0.45em}
  {\scriptsize
  这一页是在说明“一次采集窗口”在 Android 里怎么跑完：\code{EBPFCollectorService} 接收按钮、广播或实时任务入口；
  \code{DeviceCollectorDeployer} 确保手机本地有可执行的 raw eBPF collector；
  \code{memo\_libbpf\_collector} 在 root 环境 attach tracepoint 并写出 \code{MEMO\_*} 行；
  前台 app 采样器同步记录当前可见 app；
  窗口结束后停止 collector，读取 trace 文件；
  \code{EBPFTraceParser} 把文本行转成 Kotlin 事件对象，交给后面的 timeline/window 压缩。\\
  \textcolor{memoGray}{Source: \ttfamily app/src/main/java/com/memoos/realtime/RealtimeWindowRunner.kt}
  }
\end{frame}

\section{数据结构化管线}

\begin{frame}{app 序列和 eBPF 的时间对齐}
  \begin{columns}[T,totalwidth=\textwidth]
    \begin{column}{0.52\textwidth}
      \begin{block}{Foreground app sampler（前台 app 采样器）}
        每秒读取：
        \begin{itemize}
          \item \code{topResumedActivity}
          \item \code{ResumedActivity}
          \item \code{mCurrentFocus}
          \item \code{mFocusedApp}
        \end{itemize}
        然后合并成 app segment。
      \end{block}
    \end{column}
    \begin{column}{0.44\textwidth}
      \begin{block}{eBPF event alignment（内核事件时间对齐）}
        \begin{itemize}
          \item 有 wall time：直接对齐。
          \item 无 wall time：按窗口内 index 估计。
          \item 每个 timeline window 绑定 app labels 和 compressed eBPF。
        \end{itemize}
      \end{block}
    \end{column}
  \end{columns}
  \explainfoot{必须对齐是因为同一条 eBPF 证据在不同 app 场景下含义不同：Chrome 前台时大量 \code{recvfrom} 更像浏览/视频网络流，Camera 前台时文件和显示事件更像拍照链路。对齐后，MAPLE 看到的不只是“网络很多”，而是“哪个时间段、哪个 app 前台时、哪些 OS 信号一起出现”。}
\end{frame}

\begin{frame}[fragile]{timeline window 示例}
\begin{lstlisting}
{
  "index": 1,
  "start_time": "2026-07-01T22:09:23+08:00",
  "end_time": "2026-07-01T22:09:37+08:00",
  "app_labels": [
    "Clock",
    "Google Play Movies & TV"
  ],
  "app_packages": [
    "com.google.android.deskclock",
    "com.google.android.videos"
  ],
  "compressed_ebpf": [
    {"event_type": "MEMO_RECVFROM", "count": 258462, "rate_per_sec": 1367.2},
    {"event_type": "MEMO_BINDER",   "count": 134183, "rate_per_sec": 709.8}
  ]
}
\end{lstlisting}
  \explainfoot{这个 window 表示 14 秒左右的真实手机时间片：上半部分是这段时间看到的 app，下面是同一段时间压缩后的 eBPF 证据。模型可以据此判断“用户正在从 Clock 切到视频/媒体相关 app，同时网络收包和 Binder 通信很密集”。}
\end{frame}

\begin{frame}{eventType 到 category：权重矩阵}
  \scriptsize
  \begin{tabular}{lrrrrrr}
    \toprule
    eventType & Network & Binder & Runtime & File & Memory & Input \\
    \midrule
    \code{MEMO\_SENDTO} & \textbf{3} & 0 & 0 & 0 & 0 & 0 \\
    \code{MEMO\_RECVFROM} & \textbf{3} & 0 & 0 & 0 & 0 & 0 \\
    \code{MEMO\_BINDER} & 0 & \textbf{2} & 0 & 0 & 0 & 0 \\
    \code{MEMO\_SCHED} & 0 & 0 & \textbf{1} & 0 & 0 & 0 \\
    \code{MEMO\_PROCESS\_FORK} & 0 & 0 & \textbf{1} & 0 & 0 & 0 \\
    \code{MEMO\_PROCESS\_EXIT} & 0 & 0 & \textbf{1} & 0 & 0 & 0 \\
    \code{MEMO\_PROCESS\_EXEC} & 0 & 0 & \textbf{1} & 0 & 0 & 0 \\
    \code{MEMO\_OPENAT} & 0 & 0 & 0 & \textbf{1} & 0 & 0 \\
    \code{MEMO\_MEMORY} & 0 & 0 & 0 & 0 & \textbf{8} & 0 \\
    \code{MEMO\_INPUT} & 0 & 0 & 0 & 0 & 0 & \textbf{4} \\
    \midrule
    cap & 420 & 420 & 180 & 120 & 360 & 420 \\
    \bottomrule
  \end{tabular}
  \vspace{0.5em}
  \explainfoot{行是采集到的 eBPF event type，列是给 MAPLE 的语义类别。数字表示这个事件对该类别加多少分；例如一次 \code{MEMO\_RECVFROM} 给 Network 加 3 分，一次 \code{MEMO\_MEMORY} 给 Memory 加 8 分。最后一行 cap 是上限，防止网络这种高频事件把 Binder、Input、Memory 等低频但重要的信号淹没。}
\end{frame}

\begin{frame}[fragile]{Category score 如何算}
  \begin{columns}[T,totalwidth=\textwidth]
    \begin{column}{0.48\textwidth}
      \begin{block}{例：Network IO}
\begin{lstlisting}
raw counter:
  MEMO_SENDTO   count = 120
  MEMO_RECVFROM count = 180

weight:
  SENDTO -> Network +3
  RECVFROM -> Network +3

raw score = 120*3 + 180*3 = 900
cap(Network) = 420
Network IO score = min(900, 420)
\end{lstlisting}
      \end{block}
    \end{column}
    \begin{column}{0.48\textwidth}
      \begin{block}{为什么要 cap}
\begin{lstlisting}
Network events are high frequency.
If not capped:
  network may dominate all categories.

After cap:
  Network still means "very active",
  but Binder/Memory/Input can still matter.
\end{lstlisting}
      \end{block}
    \end{column}
  \end{columns}
  \vspace{0.35em}
  {\scriptsize
  Category score 是“语义强度分”：先数每类 eBPF 事件出现多少次，再乘以权重矩阵里的分数，最后用 cap 截断。
  它不是概率，也不是最终预测；它只是告诉 MAPLE 最近窗口里哪些系统活动最强。
  }
  \vspace{0.2em}
  {\scriptsize\color{memoGray}
  关键修正：Memory Management、Android Service IPC、App Process Runtime 不是 app。
  它们能说明“内存紧张 / Binder 忙 / 进程切换多”，所以影响调度动作；
  但不能显示成“下一个 app”。因此 context 会分开保存 app 候选和 OS 调度证据。
  }
\end{frame}

\begin{frame}{Context JSON 字段}
  \scriptsize
  Context JSON 是 MAPLE 每次推理真正读取的输入文件：它把“刚才用了哪些 app”和“同一时间 OS 发生了什么”压成一个结构化对象。
  \vspace{0.35em}
  \begin{tabular}{ll}
    \toprule
    Field & Meaning \\
    \midrule
    \code{historical\_app\_categories} & 用户可见 app 预测候选 \\
    \code{scheduler\_evidence\_categories} & OS/eBPF 调度证据类别 \\
    \code{historical\_app\_ids} & 类别对应的 MAPLE candidate IDs \\
    \code{prediction\_time} & 当前推理时间 \\
    \code{points\_of\_interest} & network/camera/media/display 等场景标签 \\
    \code{installed\_apps} & category 到 candidate ID 的映射 \\
    \code{app\_catalog} & 本窗口涉及的 app 元数据 \\
    \code{observed\_app\_sequence} & 带时间戳的 app segments \\
    \code{timeline\_windows} & app + compressed eBPF 对齐窗口 \\
    \code{system\_evidence} & 全局 eBPF/system 摘要文本 \\
    \code{compression} & raw rows 到 compressed signals 的元数据 \\
    \code{memory\_pressure} & normal/elevated/critical \\
    \code{available\_scheduler\_actions} & bitmask 动作表 \\
    \bottomrule
  \end{tabular}
  \explainfoot{这些字段分成三类：第一类描述“接下来可能打开什么 app”，如 app categories、app ids、installed apps；第二类描述“当前 OS 压力和证据”，如 scheduler evidence、memory pressure、system evidence；第三类描述“证据怎么来的”，如 observed app sequence、timeline windows、compression。这样模型既能预测用户行为，也能选择系统调度动作。}
\end{frame}

\begin{frame}[fragile]{真实 Context JSON 摘录}
\begin{lstlisting}[basicstyle=\ttfamily\tiny]
{
  // 用户可见 app 预测候选，不允许直接放 Binder/Memory
  "historical_app_categories": [
    "Media Codec",
    "Display Composition",
    "Network IO",
    "Communication",
    "Navigation/Location"
  ],
  // OS/eBPF 证据只作为推理和调度上下文
  "scheduler_evidence_categories": [
    "App Process Runtime",
    "Network IO",
    "Memory Management"
  ],
  "memory_pressure": "normal: no direct reclaim pressure was observed",
  // 证明模型输入来自 raw trace 的压缩摘要
  "compression": {
    "strategy": "windowed_count_aggregation",
    "raw_ebpf_rows_for_audit": 13733,
    "observed_app_sequence_rows": 6,
    "compressed_ebpf_signal_rows": 9
  }
}
\end{lstlisting}
  \explainfoot{注释强调数据边界：app 预测候选和 OS 调度证据分开。这样避免模型把 Memory/Binder 当成“下一个 app”，同时仍让这些系统信号影响 scheduler bits。}
\end{frame}

\begin{frame}{压缩不是丢证据，而是换表达}
  \begin{columns}[T,totalwidth=\textwidth]
    \begin{column}{0.52\textwidth}
      \begin{block}{Raw trace}
        \begin{itemize}
          \item 保存在 \code{/sdcard/MEMO/logs/*.trace}
          \item 用于 audit 和复现实验
          \item 可能上万行
        \end{itemize}
      \end{block}
    \end{column}
    \begin{column}{0.44\textwidth}
      \begin{block}{Model input}
        \begin{itemize}
          \item \code{event\_type}：事件类型
          \item \code{detail}：语义说明
          \item \code{count}：重复次数
          \item \code{rate\_per\_sec}：每秒频率
          \item aligned app window：对应的 app 时间窗
        \end{itemize}
      \end{block}
    \end{column}
  \end{columns}
  \vspace{0.6em}
  \centering
  \Large repeated low-level events \(\rightarrow\) compact semantic evidence
  \explainfoot{压缩的主要原因是 prompt 长度和端侧推理成本。上万行 raw eBPF 里很多行只是在重复说明同一种系统活动，逐行输入会让小模型被噪声淹没；计数后保留“发生过、发生多少、在哪个 app 时间窗发生”，这对预测和调度更有用。}
\end{frame}

\section{MAPLE 推理}

\begin{frame}{MAPLE 三阶段}
  {\scriptsize Context JSON 是 MAPLE 每次推理读取的结构化输入：里面有 app 序列、压缩 eBPF、系统状态和可选调度动作表。}
  \vspace{0.35em}
  \centering
  \resizebox{0.95\textwidth}{!}{
  \begin{tikzpicture}[
    font=\small,
    box/.style={draw, rounded corners, align=center, minimum width=3.6cm, minimum height=1.0cm},
    arrow/.style={-{Latex[length=2.2mm]}, thick}
  ]
    \node[box, fill=blue!9] (ctx) {Context JSON\\app 序列 + eBPF + 系统状态};
    \node[box, fill=orange!14, right=0.9cm of ctx] (s1) {Stage 1\\下一 app 类别};
    \node[box, fill=orange!14, right=0.9cm of s1] (s2) {Stage 2\\MAPLE app ID};
    \node[box, fill=orange!14, below=0.8cm of s2] (s3) {Stage 3\\scheduler bits 调度开关};
    \node[box, fill=green!12, left=0.9cm of s3] (out) {Android 端\\Top-3 + 系统动作};
    \draw[arrow] (ctx) -- (s1);
    \draw[arrow] (s1) -- (s2);
    \draw[arrow] (s2) -- (s3);
    \draw[arrow] (s3) -- (out);
    \draw[arrow] (out) -- (ctx);
  \end{tikzpicture}}
  \vspace{0.5em}
  \begin{block}{为什么两阶段预测}
    Stage 1 先缩小用户意图类别；Stage 2 只在该类别的 candidate IDs 内选择，降低小模型直接从复杂系统上下文生成具体 app 的难度。
  \end{block}
\end{frame}

\begin{frame}[fragile]{Stage 1 prompt：预测用户可见 app 类别}
\begin{lstlisting}
Pick the next user-facing Android app category
from real MEMO app history plus eBPF/OS context.
Answer exactly: <Category> (<percent>%).
Use only one category name from Candidates.
Do not answer Memory Management, Binder, service manager,
process runtime, or other OS-only labels as the app category.
Candidates: Media Codec, Display Composition, Network IO,
Communication, and Navigation/Location.
Memory is scheduler context, not an app category:
normal: no direct reclaim pressure was observed.
Prediction:
\end{lstlisting}
  \begin{alertblock}{意图}
    eBPF/OS 证据仍进入 prompt，但它们不再和 app category 竞争；Memory/Binder/process 是调度上下文，不是用户要打开的 app。
  \end{alertblock}
\end{frame}

\begin{frame}[fragile]{Stage 1 真实输出}
\begin{lstlisting}
[3] MODEL RAW OUTPUT:
"<Category> Media Codec (100%)"

[4] PARSED RESULT:
-> Predicted category: "Media Codec"
   confidence: 100%
\end{lstlisting}
  \explainfoot{上半部分是 MAPLE 小模型原始输出，下半部分是 Android/C++ parser 解析后的结构化结果。\code{Media Codec} 表示模型判断用户下一步更可能进入媒体播放/音视频相关 app，而不是把 \code{Memory Management}、Binder 或进程运行时误当成 app。这个结果之后会进入 Stage 2，在媒体候选 ID 中选择具体 MAPLE App ID。}
  \sourcefoot{docs/real\_device\_experiments/synthetic\_user\_30\_2026\_07\_01/\\rerun\_maple\_outputs/outputs/01\_media\_1\_scenario\_maple\_output.txt}
\end{frame}

\begin{frame}[fragile]{Stage 2 prompt：类别到候选 App ID}
\begin{lstlisting}
Choose one MAPLE candidate ID for this eBPF category.
Return one complete valid answer.
Predicted category: Media Codec.
Candidate IDs: 260.
Recent IDs: 260, 115, 245, 110, and 270.
Valid answers:
- This user will use App 260.
Prediction:

MODEL RAW OUTPUT:
"This user will use App 260."
\end{lstlisting}
  \begin{block}{说明}
    App 260 不是硬编码成某个包名；Android 端再通过本机安装应用扫描和 category 匹配映射成真实可启动 app。
  \end{block}
\end{frame}

\begin{frame}{Stage 3: scheduler\_bits 动作协议}
  \scriptsize
  \begin{tabular}{cll}
    \toprule
    bit & action id & Android action \\
    \midrule
    0 & \code{warm\_launch\_top1} & 预热 Top-1，并回 HOME \\
    1 & \code{warm\_launch\_top2\_if\_idle} & 空闲时预热 Top-2 \\
    2 & \code{trim\_memory\_low\_priority\_apps} & 对低优先级候选发 trim-memory \\
    3 & \code{kill\_selected\_background\_package} & 安全停止低优先级后台包 \\
    4 & \code{drop\_cache\_if\_critical\_memory} & 仅 critical memory 时 drop cache \\
    5 & \code{refresh\_network\_stats} & 刷新 network stats \\
    6 & \code{refresh\_service\_manager} & 刷新 service manager context \\
    7 & \code{reduce\_prewarm\_when\_display\_busy} & UI/display 忙时降低预热强度 \\
    8 & \code{skip\_camera\_prewarm\_when\_thermal\_high} & 高温时跳过 camera/media 预热 \\
    \bottomrule
  \end{tabular}
  \vspace{0.4em}
  \begin{block}{真实输出示例}
    \code{scheduler\_bits = 101010100}：执行 bit0、bit2、bit4、bit6 对应动作；具体是否落地还会经过内存/温度/前台安全检查。
  \end{block}
\end{frame}

\begin{frame}[fragile]{小模型输出解析也要工程化}
  \begin{columns}[T,totalwidth=\textwidth]
    \begin{column}{0.48\textwidth}
      \begin{block}{真实遇到的输出}
\begin{lstlisting}
"<Category>Communication (<10%))"
\end{lstlisting}
        这不是理想格式，但端侧小模型会这样输出。
      \end{block}
    \end{column}
    \begin{column}{0.48\textwidth}
      \begin{block}{parser 修正}
        \begin{itemize}
          \item 允许 \code{<Category>} prefix。
          \item 允许 \code{<10\%}。
          \item 允许多余右括号。
          \item App ID 110 在 app prediction 语境下映射为 Communication。
        \end{itemize}
      \end{block}
    \end{column}
  \end{columns}
  \vspace{0.5em}
  \begin{block}{对应代码}
    \code{llm/maple/maple\_engine/src/result\_parser.cpp}
  \end{block}
\end{frame}

\section{预测之后：真实 app 与系统调度}

\begin{frame}{从 MAPLE App ID 到真实 Top-3：映射表怎么来}
  \small
  \begin{block}{核心思想}
    MAPLE 输出的是“语义候选”（例如 Communication / App 110），不是固定包名。
    Android 端每次从本机实际安装 app 里自动找能匹配这个语义的真实应用。
  \end{block}
  \begin{enumerate}
    \item 用 \code{PackageManager}（Android 的系统应用扫描接口）扫描所有 launcher app：拿到包名、应用名、图标、launch intent。
    \item 为每个 app 自动打标签：默认浏览器/SMS/电话角色、权限、intent 能力、\code{ApplicationInfo.category}、包名和 label 关键词。
    \item 把这些标签合并成语义类别：Communication、Network IO、Camera Service、Media Codec、Payment/Security 等。
    \item 对 MAPLE 输出的类别/App ID 建候选池：只保留本机真实可启动 app，然后按匹配分排序。
  \end{enumerate}
  \vspace{0.25em}
  {\scriptsize
  排序大致是：默认角色/系统 intent 能力 \(>\) 权限和 Android category \(>\) 包名/label 关键词 \(>\) 最近使用和用户安装权重。
  所以这里没有写死“Communication = QQ”。QQ 只是这台手机上被扫描出来且匹配 Communication 的一个候选；
  换一台手机时，候选池会重新由那台手机的已安装应用生成。
  }
\end{frame}

\begin{frame}{从 MAPLE App ID 到真实 Top-3：用户看到什么}
  \begin{columns}[T,totalwidth=\textwidth]
    \begin{column}{0.48\textwidth}
      \begin{block}{同一个 MAPLE 语义在不同手机上的落点}
        \begin{itemize}
          \item 手机 A 装了 QQ：Communication \(\rightarrow\) QQ / Messages / Phone
          \item 手机 B 没装 QQ 但装了微信：Communication \(\rightarrow\) WeChat / Messages / Phone
          \item 只有系统应用：Communication \(\rightarrow\) Messages / Phone / Contacts
        \end{itemize}
      \end{block}
    \end{column}
    \begin{column}{0.48\textwidth}
      \begin{block}{用户界面的约束}
        \begin{itemize}
          \item Top-3 必须有真实包名和 launch intent。
          \item 能显示 app label 和图标。
          \item 点击后能打开对应应用。
          \item Binder、SurfaceFlinger、Memory Management 只能作为调度证据，不会显示成 app。
        \end{itemize}
      \end{block}
    \end{column}
  \end{columns}
  \explainfoot{“自然落到微信”不是模型认识微信，而是扫描器看到微信的包名、label、权限、intent 能力和通信类特征后，把它放进 Communication 候选池。MAPLE 只负责说“接下来像 Communication”，最后显示哪一个 app 由本机扫描结果决定。}
\end{frame}

\begin{frame}[fragile]{ActionExecutor 执行的是真动作}
  \begin{columns}[T,totalwidth=\textwidth]
    \begin{column}{0.49\textwidth}
\begin{lstlisting}
# warm launch Top-1
am start -W -n <component>
sleep 1
am start -a android.intent.action.MAIN \
  -c android.intent.category.HOME

# memory trim
am send-trim-memory <package> RUNNING_LOW
\end{lstlisting}
    \end{column}
    \begin{column}{0.49\textwidth}
\begin{lstlisting}
# network context refresh
dumpsys netstats

# service context refresh
service list

# critical memory only
sync; echo 1 > /proc/sys/vm/drop_caches
\end{lstlisting}
    \end{column}
  \end{columns}
  \vspace{0.4em}
  \begin{block}{安全边界}
    不杀系统包，不杀前台 app，不杀推荐候选；高内存/高温压力下减少预热而不是盲目启动更多 app。
  \end{block}
\end{frame}

\begin{frame}{实时模式任务队列}
  \centering
  \resizebox{0.95\textwidth}{!}{
  \begin{tikzpicture}[
    font=\small,
    box/.style={draw, rounded corners, align=center, minimum width=2.9cm, minimum height=0.8cm},
    arrow/.style={-{Latex[length=2mm]}, thick}
  ]
    \node[box, fill=blue!8] (w1) {Window N\\采集 app+eBPF};
    \node[box, fill=green!10, right=0.7cm of w1] (c1) {压缩\\scenario JSON};
    \node[box, fill=orange!14, right=0.7cm of c1] (m1) {MAPLE\\本地推理};
    \node[box, fill=purple!10, right=0.7cm of m1] (a1) {系统动作\\Widget 刷新};
    \node[box, fill=blue!8, below=0.8cm of w1] (w2) {Window N+1\\继续采集};
    \node[box, fill=green!10, right=0.7cm of w2] (ema) {有界\\EMA 历史记忆};
    \draw[arrow] (w1) -- (c1);
    \draw[arrow] (c1) -- (m1);
    \draw[arrow] (m1) -- (a1);
    \draw[arrow] (w1) -- (w2);
    \draw[arrow] (w2) -- (ema);
    \draw[arrow] (ema) -- (c1);
  \end{tikzpicture}}
  \vspace{0.6em}
  \begin{block}{线程控制}
    一个 collection flow + 一个 MAPLE worker；推理慢也不停止下一窗口采集，历史 memory 有长度上限，避免后台线程和上下文无限膨胀。
  \end{block}
\end{frame}

\begin{frame}{Widget：把预测变成可用的产品入口}
  \begin{columns}[T,totalwidth=\textwidth]
    \begin{column}{0.50\textwidth}
      \begin{block}{为什么需要 Widget}
        \begin{itemize}
          \item MEMO 不是让人一直打开主 App 看日志。
          \item 实时模式每个窗口更新 Top-3 后，Widget 在桌面展示真实应用图标。
          \item 用户可以直接点推荐应用，不需要理解 eBPF、Binder、scheduler bits。
        \end{itemize}
      \end{block}
    \end{column}
    \begin{column}{0.46\textwidth}
      \begin{block}{Widget 显示什么}
        \begin{itemize}
          \item Top-3 真实 app 图标和名称。
          \item 最新更新时间，说明推荐是否来自最近窗口。
          \item 点击图标直接打开对应应用。
        \end{itemize}
      \end{block}
    \end{column}
  \end{columns}
  \explainfoot{Widget 是用户层面的核心出口：eBPF 和 MAPLE 在后台工作，用户只看到“接下来可能要用的三个 app”。这让系统预测变成可点击的手机体验，而不是只停留在实验报告里。}
\end{frame}

\section{真实实验}

\begin{frame}{实验总览：我们要回答什么}
  \begin{block}{问题 1：实时产品闭环能不能连续运行？}
    对应 9 分钟实时 Top-3 变化实验。看三个连续窗口、Top-3 更新时间、MAPLE 耗时、scheduler bits。
  \end{block}
  \begin{block}{问题 2：eBPF/OS 证据是否真的帮助预测？}
    对应 30 条真实记录预测实验和理论消融。看 Top-1、Top-3 Hit、Precision/Recall/F1、MRR。
  \end{block}
  \begin{block}{问题 3：模型输出能不能进入系统调度？}
    对应 9 分钟调度记录和样本动作记录。看 warm launch、memory trim、service refresh，以及哪些动作被安全规则跳过。
  \end{block}
\end{frame}

\begin{frame}{实验之间的关系}
  \begin{itemize}
    \item 9 分钟实验验证产品形态：连续采集、连续推理、连续刷新推荐。
    \item 30 条预测实验验证预测效果：同一批样本上比较 app-only 和 app+OS/eBPF。
    \item 理论消融验证信号价值：在有明确 ground truth 的合成数据上系统删去不同证据。
    \item 调度记录验证执行闭环：MAPLE 输出不止文本，而是真的触发 Android 动作。
  \end{itemize}
  \explainfoot{这几组实验不是重复展示同一件事，而是分别回答“能不能运行”“预测有没有变好”“eBPF 是否有独立贡献”“模型能不能控制系统动作”。}
\end{frame}

\begin{frame}{9 分钟实时实验：动机与方案}
  \begin{block}{为什么做 9 分钟}
    实时模式的产品设定是每 3 分钟形成一个窗口，持续采集 app 序列、eBPF 和系统状态，然后异步 MAPLE 推理并刷新 Top-3 和调度动作。9 分钟正好覆盖 3 个完整窗口，可以观察推荐是否随场景变化。
  \end{block}
  \begin{columns}[T,totalwidth=\textwidth]
    \begin{column}{0.32\textwidth}
      \begin{block}{窗口 1}
        network/browser：浏览、网络收发、页面渲染。
      \end{block}
    \end{column}
    \begin{column}{0.32\textwidth}
      \begin{block}{窗口 2}
        camera/photo：相机、图片、显示与文件访问。
      \end{block}
    \end{column}
    \begin{column}{0.32\textwidth}
      \begin{block}{窗口 3}
        media/communication：媒体、通信、服务刷新。
      \end{block}
    \end{column}
  \end{columns}
  \explainfoot{每个窗口都在手机端真实启动已安装 app，并自然产生 eBPF 证据；不是手写 JSON，也不是 host Python 伪造 trace。Host 只负责触发和拉取结果。}
\end{frame}

\begin{frame}{9 分钟实时实验：结果}
  \scriptsize
  \begin{tabular}{lrrrl}
    \toprule
    Phase & raw eBPF & app segments & MAPLE time & Top-3 \\
    \midrule
    network & 19263 & 52 & 29106 ms & Messages / Camera / Chrome \\
    camera & 19289 & 55 & 37196 ms & QQ / Messages / Camera \\
    media/comm & 19260 & 54 & 23839 ms & Messages / Camera / QQ \\
    \bottomrule
  \end{tabular}
  \vspace{0.55em}
  \begin{block}{结果解读}
    \begin{itemize}
      \item 三个窗口都采到约 1.9 万条 raw eBPF 事件，说明采集链路持续工作。
      \item app segments 在 52-55 之间，说明实验里确实存在多次真实前台 app 变化。
      \item Top-3 在窗口间发生变化，说明推荐不是固定列表，而是随最近 app + OS 证据刷新。
    \end{itemize}
  \end{block}
  \sourcefoot{docs/real\_device\_experiments/realtime\_top3\_shift\_2026\_07\_01/}
\end{frame}

\begin{frame}[fragile]{9 分钟实验中的调度记录}
\begin{lstlisting}
phase 1 scheduler_bits = 100000000
-> maple_warm_launch_top1: Messages

phase 2 scheduler_bits = 100000000
-> maple_warm_launch_top1: QQ

phase 3 scheduler_bits = 101010100
-> maple_warm_launch_top1: Messages
-> maple_trim_memory: Camera, QQ
-> maple_drop_cache: skipped, memory not critical
-> maple_service_manager_refresh: ok
\end{lstlisting}
  \explainfoot{\code{scheduler\_bits} 是 MAPLE 输出的动作开关。第 1、2 个窗口只触发 bit0，也就是预热 Top-1。第 3 个窗口输出 \code{101010100}，同时触发 Top-1 预热、低优先级候选 memory trim、service manager refresh；drop cache 被安全规则拦截，因为当时内存不是 critical。}
  \sourcefoot{docs/real\_device\_experiments/realtime\_top3\_shift\_2026\_07\_01/}
\end{frame}

\begin{frame}{30 条真实记录预测实验：实验设计}
  \small
  \begin{columns}[T,totalwidth=\textwidth]
    \begin{column}{0.48\textwidth}
      \begin{block}{样本如何来}
        \begin{itemize}
          \item 设定 3 个 persona：media、productivity、social。
          \item Codex 根据 persona 生成 10 条短 app 使用脚本；脚本规定“已观察到的 app 序列”和“被隐藏的未来 Top-3”。
          \item Android runner 驱动真实点击/启动已安装 app。
          \item 每条样本都重新采真实 app timeline 和 raw eBPF。
        \end{itemize}
      \end{block}
    \end{column}
    \begin{column}{0.48\textwidth}
      \begin{block}{如何形成 ground truth}
        \begin{itemize}
          \item 每条样本有输入 app 序列和未来 Top-3。
          \item 未来 Top-3 在 scenario 中被 mask。
          \item MAPLE 只能看输入序列 + eBPF/OS context。
          \item 预测完成后再用隐藏 truth 评分。
        \end{itemize}
      \end{block}
    \end{column}
  \end{columns}
  \explainfoot{这里实事求是：使用意图脚本和 ground truth 是 Codex 按三类 persona 生成的；但执行阶段仍然是在 Android 设备上真实启动 app、发 input 命令、自然产生 raw eBPF。实验比较 app-only baseline 和 MEMO app+OS/eBPF，用来判断 OS 证据是否让预测更好。}
\end{frame}

\begin{frame}{30 条预测实验：评估指标}
  \small
  令预测集合为 \(P_3\)，ground truth 集合为 \(T_3\)，预测排序为 \(p_1,p_2,p_3\)。
  \vspace{0.5em}
  \begin{columns}[T,totalwidth=\textwidth]
    \begin{column}{0.50\textwidth}
      \begin{block}{集合类指标}
        \begin{itemize}
          \item Top-1 Acc: \(1[p_1 \in T_3]\)
          \item Top-3 Hit: \(1[P_3 \cap T_3 \neq \emptyset]\)
          \item Precision@3: \(|P_3 \cap T_3| / 3\)
          \item Recall@3: \(|P_3 \cap T_3| / |T_3|\)
          \item F1@3: \(2PR/(P+R)\)
        \end{itemize}
      \end{block}
    \end{column}
    \begin{column}{0.46\textwidth}
      \begin{block}{排序类指标}
        MRR: 如果第一个命中的 truth 出现在第 \(r\) 位，则得分 \(1/r\)；如果没有命中则为 0。
      \end{block}
      \begin{block}{代码如何算}
        Kotlin runner 对每条样本保存 \code{full\_memo\_top3}、\code{app\_only\_top3} 和 \code{ground\_truth\_top3}，再逐条求交集、命中位置和平均值。
      \end{block}
    \end{column}
  \end{columns}
\end{frame}

\begin{frame}[fragile]{30 条预测实验：单条指标怎么算}
\begin{lstlisting}
Predicted Top-3:
QQ / Tencent Meeting / Chrome

Ground Truth Top-3:
QQ / Chrome / Camera

Intersection size = 2  (QQ, Chrome)
Top-1 Acc = 1, because QQ is in truth
Top-3 Hit = 1, because intersection is not empty
Precision@3 = 2 / 3 = 0.667
Recall@3 = 2 / 3 = 0.667
F1@3 = 0.667
MRR = 1 / rank(QQ) = 1.0
\end{lstlisting}
  \explainfoot{这个例子有命中但不完全一致：QQ 和 Chrome 命中，Tencent Meeting 是误报，Camera 是漏报。所以 Top-3 Hit 仍然是 1，但 Precision、Recall、F1 都是 2/3。MRR 只看第一个命中的 truth 在预测列表中的位置；因为 QQ 排第一，所以 MRR 是 1。}
\end{frame}

\begin{frame}{理论消融：先在可控数据上验证 eBPF 信号}
  \begin{columns}[T,totalwidth=\textwidth]
    \begin{column}{0.52\textwidth}
      \begin{block}{实验设置}
        \begin{itemize}
          \item Chengyu Fan 的理论/合成消融实验。
          \item 1000 条带 ground truth 的合成样本。
          \item 同一个 MAPLE PromptBuilder，逐组删除 app sequence、Binder、file、memory 等证据。
          \item 目的：先回答“eBPF/OS 证据本身是否有用”。
        \end{itemize}
      \end{block}
    \end{column}
    \begin{column}{0.44\textwidth}
      \scriptsize
      \begin{tabular}{lrr}
        \toprule
        Config & DeepSeek & Qwen CPU \\
        \midrule
        Full eBPF & 88.8\% & 63.0\% \\
        Binder + File & 91.0\% & 52.8\% \\
        App sequence only & 51.1\% & 30.1\% \\
        Random & 5.6\% & 4.8\% \\
        \bottomrule
      \end{tabular}
      \vspace{0.6em}

    \end{column}
  \end{columns}
  \vspace{0.5em}
  \begin{alertblock}{和后面真实实验的一致性}
    理论实验在可控 ground truth 上显示：多源 eBPF/OS 证据显著强于仅 app 序列。
    后面的真实设备最终实验虽然样本更小、路径更接近产品运行，也得到同方向结果：
    \textbf{MEMO app+OS/eBPF} 在 Top-3 Hit、F1@3、MRR 上均超过 app-only baseline。
  \end{alertblock}
  \sourcefoot{docs/theoretical\_ablation\_experiment/synthetic\_ablation\_report.md}
\end{frame}

\begin{frame}{30 条真实记录：最终预测结果}
  \scriptsize
  \begin{tabular}{lrrrrrr}
    \toprule
    Config & Top-1 Acc & Top-3 Hit & Precision@3 & Recall@3 & F1@3 & MRR \\
    \midrule
    App-only baseline & 20.00\% & 46.67\% & 47.78\% & 47.78\% & 47.78\% & 31.11\% \\
    \textbf{MEMO app+OS/eBPF} & \textbf{23.33\%} & \textbf{50.00\%} & \textbf{52.22\%} & \textbf{52.22\%} & \textbf{52.22\%} & \textbf{34.44\%} \\
    \bottomrule
  \end{tabular}
  \vspace{0.8em}
  \begin{columns}[T,totalwidth=\textwidth]
    \begin{column}{0.48\textwidth}
      \begin{block}{相对 app-only 提升}
        \begin{itemize}
          \item Top-3 Hit: +3.33 pct points
          \item F1@3: +4.44 pct points
          \item Top-1 Accuracy: +3.33 pct points
        \end{itemize}
      \end{block}
    \end{column}
    \begin{column}{0.48\textwidth}
      \begin{block}{结论}
        full path 提升不大，但方向一致：加入 eBPF/OS context 后，Top-3 Hit、F1@3、MRR 都高于只看 app 序列。
      \end{block}
    \end{column}
  \end{columns}
  \sourcefoot{docs/real\_device\_experiments/synthetic\_user\_30\_2026\_07\_01/rerun\_inference\_only\_report.json}
\end{frame}

\begin{frame}[fragile]{单条样本的直观结果}
\begin{lstlisting}
sample_id: media_2
persona: media_browsing

MAPLE predicted category:
Communication

Predicted Top-3:
QQ / Tencent Meeting / Chrome

Ground Truth Top-3:
QQ / Tencent Meeting / Chrome

metrics:
hit_top1 = true
hit_top3 = true
F1@3 = 1.0
exact_order_top3 = true
\end{lstlisting}
  \explainfoot{这是一个 media\_browsing 风格样本：前半段脚本让设备在媒体/浏览场景中切换，真实采集里会出现网络收发、Binder、显示刷新和 app 切换证据。后半段被 mask 的未来动作是 QQ、Tencent Meeting、Chrome，语义上都和“媒体浏览后继续通信/会议/网页”相邻。Full path 结合 app 序列和 OS/eBPF 证据后把三个都排出来，所以这条样本全对。}
\end{frame}

\begin{frame}{系统调度性能实验：Stage 3 是否有用}
  \scriptsize
  \begin{block}{实验设计}
    先用真实 eBPF + MAPLE 选择同一个真实 Top-1 app：QQ。Baseline 对 QQ 做冷启动；
    Stage 3 先执行同一组 \code{scheduler\_bits=101010100}，再启动同一个 QQ。
    因此比较的是“同一目标 app 无调度”与“同一目标 app 经 MAPLE 调度后”的系统指标。
  \end{block}
  \begin{tabular}{lrrr}
    \toprule
    Metric & No scheduler & Stage 3 & Improvement \\
    \midrule
    Launch TotalTime & 433 ms & 168 ms & 61.20\% \\
    WaitTime & 440 ms & 173 ms & 60.68\% \\
    综合压力分数 & 422.56 & 9.71 & 97.70\% \\
    CPU busy & 26.70\% & 20.79\% & 22.13\% \\
    iowait & 0.094\% & 0.023\% & 75.42\% \\
    reclaim & 10046 & 0 & 100.00\% \\
    jank rate & 1.240\% & 0.886\% & 28.56\% \\
    \bottomrule
  \end{tabular}
  \vspace{0.45em}
  \explainfoot{这页只评价系统调度性能，不评价预测准确率。正数表示 lower-is-better 指标下降。Stage 3 真实执行了 7 条动作记录，其中包括 Top-1 warm launch、memory trim、service manager refresh；drop cache 被安全规则跳过，因为当时内存不是 critical。}
  \sourcefoot{docs/real\_device\_experiments/system\_scheduler\_performance/latest\_pressure\_experiment.json}
\end{frame}

\section{实现文件与结论}

\begin{frame}{关键代码：三分钟采集窗口}
  \begin{block}{\code{RealtimeWindowRunner.kt}}
    这个文件负责在手机内完成一个窗口的真实采集：启动 raw eBPF collector，同时采样前台 app，最后输出 app timeline、raw trace 和 timeline windows。
  \end{block}
  \begin{block}{为什么关键}
    它决定了“输入是不是来自真实使用”。如果这里仍然依赖 host Python 或手写 JSON，后面的 MAPLE 推理和调度都不能算端侧产品闭环。
  \end{block}
  \sourcefoot{app/src/main/java/com/memoos/realtime/RealtimeWindowRunner.kt}
\end{frame}

\begin{frame}{关键代码：有界实时记忆}
  \begin{block}{\code{RealtimeMemoryStore.kt}}
    这个文件维护历史 app 类别和资源压力的 EMA，也就是指数衰减平均。它只保留固定长度的历史窗口，不让后台 memory 越跑越大。
  \end{block}
  \begin{block}{为什么关键}
    实时模式不能只看最近三分钟，也不能无限保存所有历史。EMA 让模型知道用户近期偏好，同时控制端侧开销。
  \end{block}
  \sourcefoot{app/src/main/java/com/memoos/realtime/RealtimeMemoryStore.kt}
\end{frame}

\begin{frame}{关键代码：MAPLE 输入构造}
  \begin{block}{\code{MapleScenarioBuilder.kt}}
    这个文件把 app 序列、eBPF counter、系统状态、安装应用表、scheduler 动作表组合成 MAPLE 使用的 scenario JSON。
  \end{block}
  \begin{block}{为什么关键}
    它是“app + OS 到 app + OS”的数据边界：既要把 app 预测候选给模型，也要把 OS 证据作为调度上下文给模型。
  \end{block}
  \sourcefoot{app/src/main/java/com/memoos/maple/MapleScenarioBuilder.kt}
\end{frame}

\begin{frame}{关键代码：本地模型调用}
  \begin{block}{\code{MapleShellBackend.kt}}
    这个文件在手机本地调用 MAPLE/llama.cpp 可执行文件，传入 scenario JSON，拿回 Top-3 预测、pressure analysis 和 scheduler bits。
  \end{block}
  \begin{block}{为什么关键}
    它保证预测不是电脑上跑的脚本，而是 Android 设备侧本地推理；拔掉数据线后仍然可以通过 App/Widget 运行。
  \end{block}
  \sourcefoot{app/src/main/java/com/memoos/maple/MapleShellBackend.kt}
\end{frame}

\begin{frame}{关键代码：真实应用映射}
  \begin{block}{\code{AppIdMapping.kt}}
    这个文件扫描本机安装 app，把 MAPLE 的类别/App ID 映射成可启动的 QQ、Chrome、Camera 等真实应用。
  \end{block}
  \begin{block}{为什么关键}
    用户看到的是应用，不是进程、Binder 或 SurfaceFlinger。这个映射层把模型输出变成用户能理解、能点击的 Top-3。
  \end{block}
  \sourcefoot{app/src/main/java/com/memoos/action/AppIdMapping.kt}
\end{frame}

\begin{frame}{关键代码：系统调度动作}
  \begin{block}{\code{ActionExecutor.kt}}
    这个文件解释 scheduler bits，执行 warm launch、memory trim、network/service refresh，并做内存和温度安全检查。
  \end{block}
  \begin{block}{为什么关键}
    MAPLE 输出不能只停留在“建议”。ActionExecutor 把模型 bitmask 变成可执行 Android 动作，同时避免杀前台 app、系统包或推荐候选。
  \end{block}
  \sourcefoot{app/src/main/java/com/memoos/action/ActionExecutor.kt}
\end{frame}

\begin{frame}{关键代码：桌面 Widget}
  \begin{block}{\code{MemoWidgetProvider.kt}}
    这个文件在桌面展示 Top-3 真实应用和最新更新时间，点击图标可以直接打开对应应用。
  \end{block}
  \begin{block}{为什么关键}
    Widget 是用户真正会用到的入口。后台的 eBPF、MAPLE 和调度动作不需要暴露给用户；用户只需要看到“接下来可能要打开的三个 app”。
  \end{block}
  \sourcefoot{app/src/main/java/com/memoos/widget/MemoWidgetProvider.kt}
\end{frame}

\begin{frame}{关键代码：提示词与输出解析}
  \begin{block}{\code{prompt\_builder.cpp} / \code{result\_parser.cpp}}
    这两个文件约束小模型输出 app 类别、App ID 和 scheduler bits，并解析不规整的模型文本。
  \end{block}
  \begin{block}{为什么关键}
    端侧小模型不够“听话”，输出可能带多余括号或解释文字。parser 必须把这些脏输出尽量转成结构化结果，才能驱动推荐和调度。
  \end{block}
  \sourcefoot{llm/maple/maple\_engine/src/prompt\_builder.cpp and llm/maple/maple\_engine/src/result\_parser.cpp}
\end{frame}

\begin{frame}{最终系统做到了什么}
  \begin{block}{端侧闭环}
    \centering
    手机本地 app 序列 + raw eBPF + 系统状态\\
    \(\Downarrow\)\\
    压缩后的 scenario JSON（MAPLE 输入）\\
    \(\Downarrow\)\\
    本地 MAPLE：Top-3 真实 app + scheduler bits 调度开关\\
    \(\Downarrow\)\\
    Android Widget 展示 + ActionExecutor 系统调度
  \end{block}
  \begin{itemize}
    \item app+OS 共同进入模型，不再只是浅层 app sequence（应用序列）。
    \item 模型输出不止推荐，也控制预定义系统动作。
    \item Widget 让 Top-3 推荐成为桌面可点击入口，方便用户直接打开预测应用。
    \item 真实 Pixel 5 上完成 eBPF 采集、推理、调度、Widget 展示和实验归档。
  \end{itemize}
\end{frame}

\begin{frame}{限制与下一步}
  \begin{columns}[T,totalwidth=\textwidth]
    \begin{column}{0.49\textwidth}
      \begin{block}{当前限制}
        \begin{itemize}
          \item 小模型能力有限，不够“听话”：即使 prompt 要求固定格式，它也可能多输出括号、解释或错误字段。
          \item 目前已经有 parser 处理脏输出，但还不够 robust（鲁棒）：需要覆盖更多奇怪输出。
          \item 实时性的代价是后台常驻：系统持续采集，每三分钟推理一次；实时性和轻量性是一对 tradeoff。
        \end{itemize}
      \end{block}
    \end{column}
    \begin{column}{0.49\textwidth}
      \begin{block}{下一步}
        \begin{itemize}
          \item 更细粒度 scheduler policy：把 bitmask 从粗粒度动作扩展成带强度、时机和安全阈值的调度策略。
          \item 强化 app mapping：更自动地识别新安装 app 的语义类别，减少依赖包名关键词。
          \item 完善 capability report 和部署脚本：让设备端清楚报告 root、BTF、tracefs、模型文件和 collector 状态。
        \end{itemize}
      \end{block}
    \end{column}
  \end{columns}
\end{frame}

\begin{frame}
  \centering
  \Huge Thanks
  \vspace{1em}

  \large MEMO-Appflow: app + OS evidence \(\rightarrow\) local AI reasoning \(\rightarrow\) app + OS actions
\end{frame}

\end{document}
